\documentclass[10pt,a4paper]{amsart}
\usepackage[margin=19mm]{geometry}
\usepackage{amsmath,amssymb,mathtools}
\usepackage[scr=rsfs,cal=euler]{mathalfa}
\usepackage{xfrac}
\usepackage[unicode,hidelinks,bookmarks=false]{hyperref}
\newcommand{\ie}{i.e.\ }
\newcommand{\cf}{cf.\ }
\newcommand{\resp}{respectively\ }
\newcommand{\Subsection}{\subsection}
\providecommand{\nobreakdash}{\nobreak\hbox{-}}
\setlength{\emergencystretch}{2em}
\multlinegap0pt
\usepackage{bm}
\usepackage{bbm}
\usepackage{dsfont}
\usepackage{xspace}
\usepackage{cancel}
\usepackage{mathtools}
\usepackage{amsthm}
\makeatletter
\@ifundefined{lemma}{\newtheorem{lemma}{Lemma}}{}
\makeatother
\title[Valuation Measure of the Stock Market using Stochastic Volat]{Valuation Measure of the Stock Market using Stochastic Volatility and Stock Earnings}
\author{Andrey Sarantsev, Angel Piotrowski, Ian Anderson}
\date{Source: arXiv:2508.06010v2 (CC BY 4.0)}
\begin{document}
\maketitle

\noindent\emph{Source and licence.} Valuation Measure of the Stock Market using Stochastic Volatility and Stock Earnings. Version arXiv:2508.06010v2 (CC BY 4.0), distributed under the Creative Commons Attribution 4.0 International licence, as recorded in the arXiv licence metadata for this exact version and shown on the abstract page https://arxiv.org/abs/2508.06010v2. Full rights notice: licenses/arxiv-2508.06010-CC-BY-4.0.txt.\par\smallskip

\noindent\textit{Editorial scope.} Counted unit: the Lemma whose source label is \texttt{eq:finite-phi} in the article (source0.tex lines 752--762, source label \texttt{eq:finite-phi}), delivered together with the article's own complete original proof of it (source0.tex lines 764--814, one \texttt{proof} environment of 51 lines). No mathematics is written, repaired or re-derived: statement and proof are the article's own text line for line. Required context retained uncounted: none. Every definition, display and label of those lines is retained unchanged. Omitted from this package and not redistributed: 1--751, 763--763, 815--998. Nothing inside a retained excerpt is omitted or altered: every retained body is delivered exactly as the article's lines are written, so no omission receipt of any kind accompanies this package. Citations: the retained excerpts carry no citation command at all, so no bibliography entry of the article is transcribed and no reference is redistributed. Reference adaptation: none was needed and none was made; every reference inside the delivered unit resolves inside it, so no unresolved reference or citation is produced. Printed slips kept literal: no printed word, symbol, hypothesis or number of the counted statement or of its proof was altered. Layout and class: the article's own class is \texttt{amsart} with packages including geometry, graphicx, subfig; the wrapper is the lane's pinned \texttt{amsart} editorial wrapper at 10pt on A4, which restates the theorem-like environments the retained text needs and adds the article's own macro definitions that the retained text needs (none), with extra wrapper packages (bm, bbm, dsfont, xspace, cancel, mathtools). The delivered numbering is the unit's own. Assets: no retained span contains a figure environment, a graphics-inclusion command, a PDF-page inclusion, a table or a TikZ picture; no image file is copied into this package, the package declares no asset dependency and no image file is redistributed with it. Licence: version arXiv:2508.06010v2 of this article is distributed under the Creative Commons Attribution 4.0 International licence, as recorded in the arXiv licence metadata for this exact version and shown on the abstract page https://arxiv.org/abs/2508.06010v2. A full rights notice is delivered at licenses/arxiv-2508.06010-CC-BY-4.0.txt. Characters: every retained excerpt is pure ASCII, so the delivered engine, XeLaTeX with its default Latin Modern text font, reports no missing character.
\begin{lemma} If two random variables $X, Y$ (not necessarily independent) have 
\begin{equation}
\label{eq:finite-phi}
\mathbb E\varphi(|X|) < \infty,\quad \mathbb E\varphi(|Y|) < \infty,
\end{equation}
then the sum and product of these two random variables also satisfy these properties:
$$
\mathbb E\varphi(|X + Y|) < \infty,\quad \mathbb E\varphi(|XY|) < \infty.
$$
\label{lemma:sum-product}
\end{lemma}

\begin{proof} First, let us show this for $X + Y$. This function $\varphi$ is nondecreasing but not concave:
$$
\varphi(x) = 
\begin{cases} 
\ln x, \, x \ge 1;\\ 
0,\, x \in [0, 1]. 
\end{cases}
$$
But it has a {\it concave hull:} the smallest concave $\psi : \mathbb R_+ \to \mathbb R$ such that $\psi(x) \ge \varphi(x)$ for all $x \ge 0$. We can find it by drawing a tangent line from $(0, 0)$ to $x \mapsto \psi(x) = \ln x$ at the point $(a, \ln a)$ for $a > 1$. We compare the slope of this line: $\ln a/a = \psi'(a) = 1/a$, equivalent to  $\ln a = 1$ and $a = e$. Therefore,
$$
\varphi(x) \le \psi(x) = 
\begin{cases}
x/e,\, x \le e;\\
\ln x,\, x \ge e.
\end{cases}
$$
This function $\psi$ satisfies 
\begin{equation}
\label{eq:psi-linearity}
\psi(x + y) \le \psi(x) + \psi(y),\quad x, y \ge 0.
\end{equation}
Indeed, from concavity of $\psi$, we have: $\psi(x + y) - \psi(x) \le \psi(y) - \psi(0) = \psi(y)$. Thus for any random variable $X$, 
\begin{equation}
\label{eq:phi-psi}
\mathbb E\varphi(|X|) < \infty \Leftrightarrow \mathbb E\psi(|X|) < \infty.
\end{equation}
Applying~\eqref{eq:phi-psi} to~\eqref{eq:finite-phi}, we get: 
\begin{equation}
\label{eq:finite-psi}
\mathbb E\psi(|X|) < \infty,\quad \mathbb E\psi(|Y|) < \infty.
\end{equation}
Applying~\eqref{eq:psi-linearity} to~\eqref{eq:finite-psi}, we get: $\mathbb E\psi(|X + Y|) < \infty$. Finally, apply~\eqref{eq:phi-psi} again and complete the proof that~\eqref{eq:finite-phi} implies $\mathbb E\varphi(|X + Y|) < \infty$. Next, we prove this for $XY$. Indeed, 
\begin{equation}
\label{eq:varphi-product}
\varphi(xy) \le \varphi(x) + \varphi(y),\quad x, y \ge 0.
\end{equation}
To prove~\eqref{eq:varphi-product}, we need to consider four cases, comparing $x, y, xy$ with $1$: 
\begin{itemize}
\item if $x, y \le 1$, then $xy \le 1$, so $\varphi(x) = \varphi(y) = \varphi(xy) = 0$;
\item if $x, y \ge 1$, then $xy \ge 1$, so $\varphi(x) = \ln x$, $\varphi(y) = \ln y$, and $\varphi(xy) = \ln(xy) = \ln x + \ln y$;
\item if $x \le 1$ and $y \ge 1$ so that $xy \le 1$, then $\varphi(x) = 0$ and $\varphi(y) = \ln y$, but 
$$
\varphi(xy) = 0 \le 0 + \ln y = \varphi(x) + \varphi(y); 
$$
\item if $x \le 1$ and $y \ge 1$ so that $xy \ge 1$, then $\varphi(x) = 0$ and $\varphi(y) = \ln y$, but 
$$
\varphi(xy) = \ln(xy) = \ln x + \ln y \le 0 + \ln y = \varphi(x) + \varphi(y).
$$
\end{itemize}
This completes the proof of~\eqref{eq:varphi-product}, and with it the lemma.
\end{proof}

\end{document}
